% Generated by render_results.py; do not hand-edit.
The build uses \code{leanprover/lean4:v4.19.0}, with Mathlib at \nolinkurl{c44e0c8ee63ca166450922a373c7409c5d26b00b}. The audit counts source theorem declarations separately from generated logical helpers. It records 14 compiler runtime extraction auxiliaries outside the proof inventory; these are not used to replace kernel checking. The reported counts and names come from the build audit JSON.

\begin{center}

\begin{tabular}{p{0.29\textwidth}p{0.65\textwidth}}

\toprule

Claim family & Compiled declarations in namespace \code{StateCut} \\

\midrule

Moment and residual bounds & \code{abs\_chord\_sum}, \code{moment\_residual\_sound}, \code{finite\_abs\_sum\_bound}, \code{finite\_residual\_bound} \\

Finite-count strength and feasibility & \code{finite\_envelope\_le\_chord}, \code{finite\_envelope\_attained}, \code{floor\_remainder\_contract} \\

Rounding-cell composition & \code{residual\_to\_cell}, \code{moment\_to\_cell}, \code{finite\_to\_cell} \\

Conditional numerical bridge & \code{backend\_bridge\_gate}, \code{backend\_bridge\_cell} \\

Disjoint frontier sums & \code{disjoint\_split\_sum}, \code{frontier\_sum\_bounds} \\

Persistent writes and traces & \code{singleton\_writes}, \code{checked\_write\_frontier}, \code{write\_frontier\_future} \\

\bottomrule

\end{tabular}

\end{center}

The finite-count theorem takes an integer endpoint count and a remainder satisfying the summary equation. The separate mathematical floor/remainder theorem supplies those conditions in the nonconstant, non-upper-endpoint case. Constant cases are proved separately. Neither result establishes that Python's runtime executes the specified mathematics; the exact-rational tests check that implementation separately. The cell theorems take a correct rounding-cell contract as a premise and do not prove the full BF16 encoding/cell generator. The frontier results concern disjoint sums; they do not formally verify the Python persistent forest or its memory-access counters.
